%% Liaison complète tige-piston, en demi-coupe : deux surfaces fonctionnelles, l'une cylindrique
%% (positionnement radial), l'autre plane (positionnement axial), plus la vis de maintien.
%% Le rôle de chaque surface est la réponse : il passe par \rappel.
\documentclass[tikz,border=4pt]{standalone}
\usepackage{liaisons}
\begin{document}
\begin{tikzpicture}[x=1cm,y=1cm,
  axe/.style={line width=0.5pt, draw=black!55, dash pattern=on 9pt off 3pt on 2pt off 3pt},
  repere/.style={line width=0.5pt, draw=black!60},
  etiq/.style={font=\small, text=black!80, inner sep=2pt}]

%% ---------------- la tige (1), demi-coupe au-dessus de l'axe ----------------
\hachures[draw=solideA, line width=0.9pt]{
  (0,0) -- (0,1.05) -- (2.4,1.05) -- (2.4,0.62) -- (4.6,0.62) -- (4.6,0) -- cycle}
\node[font=\small\bfseries, text=solideA, anchor=south] at (1.2,1.12) {Tige (1)};

%% ---------------- le piston (2), enfilé sur la portée cylindrique ----------------
\hachuresb[draw=solideB, line width=0.9pt]{
  (2.4,0.62) -- (2.4,1.95) -- (3.9,1.95) -- (3.9,0.62) -- cycle}
\node[font=\small\bfseries, text=solideB, anchor=south] at (3.15,2.05) {Piston (2)};

%% ---------------- la vis de maintien ----------------
\draw[line width=0.8pt, draw=black!70, fill=black!14] (4.6,0) rectangle (5.2,0.45);
\draw[line width=0.8pt, draw=black!70] (4.05,0) rectangle (4.6,0.2);
\node[etiq, anchor=west, text=black!65] at (5.28,0.25) {vis de maintien};

%% ---------------- l'axe de révolution ----------------
\draw[axe] (-0.55,0) -- (6.0,0);
\node[etiq, anchor=east, text=black!55, font=\scriptsize] at (-0.58,0) {axe};

%% ---------------- SF1 : la portée cylindrique ----------------
\draw[line width=2.4pt, draw=solideE, line cap=round] (2.45,0.62) -- (3.85,0.62);
\draw[repere] (3.5,0.62) -- (4.35,-0.72);
\node[etiq, text=solideE, anchor=north west, font=\small\bfseries] at (4.1,-0.68) {SF1 : surface cylindrique};
\node[etiq, anchor=north west, align=left, text=black!70] at (4.1,-1.05)
  {\rappel{positionnement \textbf{radial}\\ (la coaxialité)}};

%% ---------------- SF2 : la face plane de l'épaulement ----------------
\draw[line width=2.4pt, draw=solideF, line cap=round] (2.4,0.66) -- (2.4,1.91);
\draw[repere] (2.4,1.5) -- (1.35,2.35);
\node[etiq, text=solideF, anchor=south east, font=\small\bfseries] at (1.4,2.35) {SF2 : surface plane};
\node[etiq, anchor=north east, align=right, text=black!70] at (1.4,2.32)
  {\rappel{positionnement \textbf{axial}\\ (la butée)}};
\end{tikzpicture}
\end{document}
